% Complete proof candidate. Individual check status is in result_audit.json.
\subsection{The precise nonarchimedean comparison}
\begin{proposition}[ARI C.7.8; Known with derivation: a nonzero ordinary height]\label{ari:C.7.8}
Let $E/\Q$ have complex multiplication by the full ring of integers
of an imaginary quadratic field $K$ of class number one. Assume
that its functional-equation sign is $-1$, and let $p>3$ be a
prime of good ordinary reduction, $p\mathcal O_K=\mathfrak p
\mathfrak p^*$. For a nontorsion $P\in E(\Q)$ the
$\mathfrak p$-adic height and the cyclotomic $p$-adic height
of $P$ over $K$, in the normalization of
\citet[Appendix A]{agboola2006}, are nonzero.
In particular this applies to
\[
 E:y^2+y=x^3-38x+90,\qquad P=(0,9),\qquad p=5.
\]
\end{proposition}
\begin{proof}
We use the exact theorem reproduced as Theorem A.1 in Rubin's
appendix to \citet{agboola2006}, attributed there to
\citet{bertrand1984}: for a number field $F$ and a point of
infinite order in $E(F)$, its height $h_{F,\mathfrak p}$ is
nonzero. We retain the standing hypotheses of that appendix,
including the sign assumption, rather than infer a wider theorem
from the abbreviated statement. The height takes values in
the one-dimensional $\Q_p$-space attached to the unique
$\Z_p$-extension unramified outside $\mathfrak p$.
Taking $F=K$ proves its nonvanishing. The polarized height of
a rational point has zero anticyclotomic projection: complex
conjugation fixes the point and acts as $-1$ on that projection.
Lemma A.2 of the same appendix states that if two of the
cyclotomic, anticyclotomic and $\mathfrak p$ projections
vanish, then the third vanishes. It follows that the cyclotomic
projection is nonzero as well. This argument asserts nonvanishing
on this point, not nondegeneracy of a height pairing on an
arbitrary Mordell--Weil space tensored with $\Q_p$.

For the displayed curve, C.6.R1--R2 prove nontorsion and
analytic rank one, hence sign $-1$. Its CM order has
discriminant $-19$, which is maximal, and its class number is
one: the reduced positive binary forms of that discriminant
give only $(1,1,5)$ up to the usual boundary convention.
The discriminant of the minimal curve is $-19^3$, so reduction
at $5$ is good. Counting directly over $\F_5$ gives seven
points, hence $a_5=-1$ and ordinary reduction; equivalently
$-19$ is a nonzero square modulo $5$. Every hypothesis of the
cited theorem therefore holds in this concrete unresolved
archimedean example.
\end{proof}

The comparison explains the missing step in an attempted
archimedean adaptation. The ordinary $p$-adic sigma construction
uses a holomorphic quadratic correction. In the CM normalization
the corresponding complex holomorphic function is $\theta$
of C.7.6, whereas the real metric also contains
$\exp(-\pi|u|^2/(2A))$. Theorem A.1 supplies no identity
between the ordinary $p$-adic height and that real factor.
Transferring the conclusion would require precisely such a
new argument; replacing the real factor by the holomorphic
CM correction does not preserve the norm being asked about.

There is also a definite limitation on the maximal-compact
approach. \citet[Theorem 2.1]{bertrand1995} constructs
nontrivial multiplicative extensions of polarized abelian
varieties with nontorsion algebraic points in their maximal
real compact subgroup. Its hypotheses include an endomorphism
not fixed by the Rosati involution and a nontorsion rational
point over a real number field. The construction uses a base
point related to the extension class by a Rosati-antisymmetric
endomorphism. It cannot be replaced by a theorem saying that
every algebraic compact point of every multiplicative
extension is torsion. For our diagonal class, an equality
$NP=f(P)$ with $N>0$ and $f+f^\dagger=0$ would force
$P$ torsion: $N-f$ has no zero eigenvalue, since
$(N-f)^\dagger(N-f)=N^2-f^2$ is positive. In the elliptic
case this is simply a nonzero isogeny. The diagonal positivity
used in C.7.3--C.7.4 is therefore essential, but these
counterexamples neither prove nor refute the desired
unconditional diagonal statement. The additive result in
\citet[Theorem 3.1]{bertrand1995}, and its higher-dimensional
extension in \citet{bost2009}, requires a point in a vector
extension, a hypothesis not furnished by C.7.2.
